\inicializarSeccion{Discusión y Conclusiónes}
A partir de nuestra interpretación, podemos confirmar el enunciado inicial de la situación problema: El océano cambia, los datos lo revelan. Pero estos cambios cómo pueden estar para bien, también pueden estar para mal. Es innegable que con la transición ambiental ocurrida en cada uno de los ambientes ocurre un desplazamiento parcial o total de la flora y fauna marina. Debido a esto último, no nos debe de extrañar que especies que antes dominaban ciertas zonas; o bien, áreas en donde antes abundaba la vida marina; dejarán de hacerlo por estos cambios. Es importante recalcar que esto no solo afecta a la vida marina; los datos son claros; los patrones mostrados en el presente estudio perjudican a las actividades humanas como la pesca, la exploración marítima, el turismo, entre otros….

Con nuestras figuras, pudimos visualizar que dependiendo del ambiente, ciertas zonas experimentan ciertas condiciones ambientales muy variadas; inusualmente variadas. Hay algunas variables como los fosfatos en la figura [REF FIGURA](4.3.6) que contienen varianzas atípicas incluso extremas. Pero también tenemos casos que son inusualmente estables. Tal como es el caso de la figura [REF FIGURA](4.3.2) en donde aparenta haber una gran estabilidad en temperatura común desde “very shallow” hasta “intermediate”. O bien, en la figura [REF FIGURA](4.3.3) en donde la diferencia entre el número de fosfatos y que tan densa es el agua contiene una mediana muy similar. Estos valores son peculiares porque uno esperaría que decreciera mientras más profundidad (y por ende más densidad) el agua contiene. 

Dentro este análisis y reflexión, nosotros queremos dar enfoque en tres preguntas que definirán las siguientes etapas de la situación problema.

\begin{enumerate}
    \item \textit{¿De qué manera las variaciones e incrementos alarmantes en la temperatura y la salinidad del agua en las últimas décadas ponen en riesgo el equilibrio del ecosistema y la biodiversidad marina?}
    \item \textit{¿Qué impacto tienen los valores atípicos y las condiciones ambientales inusualmente variadas de clorofila-a y fosfatos en las actividades humanas como la pesca y la exploración marítima?}
    \item \textit{¿En qué medida la disminución del oxígeno disuelto y de la clorofila-a en las zonas con mayor densidad de agua funciona como un indicador del deterioro o transformación drástica de un ecosistema marino específico?}
\end{enumerate}

Usando estas preguntas, queremos encontrar respuestas acerca del origen de los datos. No necesariamente por una falta de validez, si no para ver cómo estas variables locales afectan a todo el ecosistema global; indagando no sólo en la rigidez de los datos, sino en plantear un origen de la problemática y con ello una posible solución.




